\references{
\frfilename{mt1r.tex}
\versiondate{31.5.03}

\gdef\topparagraph{}
\gdef\bottomparagraph{}
\gdef\newparagraph{}
\gdef\chaptername{Referensi}
\gdef\sectionname{Referensi}
\gdef\headlinesectionname{Referensi}

\centerline{\bf Referensi untuk Jilid 1}

\medskip

Selain karya-karya (yang sangat sedikit) yang telah saya sebutkan dalam
jilid ini, saya mencantumkan beberapa buku yang saya sendiri gunakan untuk
mempelajari teori ukuran, sebagai ungkapan rasa hormat dan terima kasih,
serta untuk memberi Anda kesempatan mencoba berbagai pendekatan lain.

\medskip

Bartle R.G.\ [66] {\it The Elements of Integration.}  Wiley, 1966.

Cohn D.L.\ [80] {\it Measure Theory.}  Birkh\"auser, 1980.

Facenda J.A. \& Freniche F.J.\ [02]
{\it Integraci\'on de funciones de varias variables.}
Ediciones Pir\'amide, 2002.

Halmos P.R.\ [50] {\it Measure Theory.}   Van Nostrand, 1950.

Lebesgue H.\ [1901] `Sur une g\'en\'eralisation de l'int\'egrale
d\'efinie', C.R.Acad.Sci.\ (Paris) 132 (1901) 1025-1027.
Dicetak ulang dalam {\smc Lebesgue 72}, jil.\ 1.

\leaveitout{Lebesgue H.\ [1902] {\it Int\'egrale, Longueur, Aire,}  th\`ese pr\'esent\'ee \`a la Facult\'e de Sciences de Paris,  Annali di Matematica Pura Appl.\ 7 (1902).   Dicetak ulang dalam {\smc
Lebesgue 72}, jil.\ 1.}

Lebesgue H.\ [1904] {\it Le\c{c}ons sur l'int\'egration et la recherche
des fonctions primitives.} Gauthier-Villars, 1904.
Dicetak ulang dalam {\smc Lebesgue 72}, jil.\ 2.
\cmmnt{[Jil.\ 1 {\it pendahuluan.}]}

Lebesgue H.\ [1918] `Remarques sur les th\'eories de la mesure et de l'int\'egration', Ann.\ Sci.\ Ecole Norm.\ Sup.\ 35 (1918) 191-250.
Dicetak ulang dalam {\smc Lebesgue 72}, jil.\ 2.
\cmmnt{[Jil.\ 1 {\it pendahuluan.}]}

Lebesgue H.\ [72] {\it Oeuvres Scientifiques.}   L'Enseignement Math\'ematique, Institut de Math\'ematiques, Univ.\ de Gen\`eve, 1972.

Munroe M.E.\ [53] {\it Introduction to Measure and Integration.}
Addison-Wesley, 1953.

Royden H.L.\ [63] {\it Real Analysis.}  Macmillan, 1963.

Solovay R.M.\ [70] `A model of set theory in which every set of reals
is Lebesgue measurable', Annals of Math.\ 92 (1970) 1-56.   \cmmnt{[134C.]}

Widom H.\ [69] {\it Lectures on Measure and Integration.}  Van Nostrand Reinhold, 1969.

Williamson J.H.\ [62] {\it Lebesgue Integration.}  Holt, Rinehart \& Winston, 1962.

\frnewpage

}%end of references
